
\vspace{1cm}

\noindent  \textbf{Leia as instruções antes de fazer a lista:}

\vspace{0.3cm}
\noindent A resolução deve ser entregue em formato Jupyter Notebook ou Google Colab na plataforma Google Classroom, na tarefa designada e até o dia determinado.

\noindent  Não é necessário testar se os dados passados por argumento são válidos a menos que pedido na questão.


\vspace{0.5cm}


\begin{enumerate}[label=\textbf{Q\arabic*.}]



    % Questão 1
    \item Escreva uma função em Python que busca todas as ocorrências de um determinado número em uma lista.
    \begin{enumerate}[label=\alph*)]
        \item A função deve receber por argumento uma lista de números e um número (nessa ordem), e retornar uma nova lista contendo as posições (índices) em que esse número aparece na lista passada por argumento. Se o número não existir na lista, a função deve retornar uma lista vazia.
        \item Teste a sua função de acordo com as entradas e saídas esperadas a seguir:
        \begin{itemize}
            \item \textbf{Entrada:} \code{busca\_posicoes([10, 20, 30, 20, 40, 20], 20)} $\rightarrow$ \textbf{Saída Esperada:} \code{[1, 3, 5]}
            \item \textbf{Entrada:} \code{busca\_posicoes([1, 2, 3, 4, 5], 99)} $\rightarrow$ \textbf{Saída Esperada:} \code{[]}
            \item \textbf{Entrada:} \code{busca\_posicoes([], 5)} $\rightarrow$ \textbf{Saída Esperada:} \code{[]}
        \end{itemize}
    \end{enumerate}

    \vspace{0.5cm}

    \ifmostrarGabarito
\begin{minted}{python}
# Resposta Questão 1:
def busca_posicoes(lista, numero):
    posicoes = []
    i = 0
    while i < len(lista):
        if lista[i] == numero:
            posicoes.append(i)
        i += 1
    return posicoes

# Testes (Item b):
print(busca_posicoes([10, 20, 30, 20, 40, 20], 20)) # Esperado: [1, 3, 5]
print(busca_posicoes([1, 2, 3, 4, 5], 99))            # Esperado: []
print(busca_posicoes([], 5))                         # Esperado: []
\end{minted}
\fi


    % Questão 2
    \item O produto escalar entre dois vetores $a$ e $b$ de tamanho $n$ é definido por:
    \[
    [a_{1}, a_{2}, \dots, a_{n}] \cdot 
    \begin{bmatrix}
    b_{1} \\
    b_{2} \\
    \vdots \\
    b_{n}
    \end{bmatrix}
    = a_{1} \cdot b_{1} + a_{2} \cdot b_{2} + \dots + a_{n} \cdot b_{n}
    \]
    \begin{enumerate}[label=\alph*)]
        \item Escreva uma função em Python que recebe por argumento duas listas que representam vetores e retorna o resultado do produto escalar entre elas. Considere que as listas recebidas possuem o mesmo tamanho e que todos os seus elementos são números.
        \item Teste a sua função de acordo com as entradas e saídas esperadas a seguir:
        \begin{itemize}
            \item \textbf{Entrada:} \code{produto\_escalar([1, 2, 3], [4, 5, 6])} $\rightarrow$ \textbf{Saída Esperada:} \code{32}
            \item \textbf{Entrada:} \code{produto\_escalar([2.5, 0.0], [4.0, 10.0])} $\rightarrow$ \textbf{Saída Esperada:} \code{10.0}
            \item \textbf{Entrada:} \code{produto\_escalar([], [])} $\rightarrow$ \textbf{Saída Esperada:} \code{0}
        \end{itemize}
    \end{enumerate}

    \vspace{0.5cm}
    \ifmostrarGabarito{
\begin{minted}{python}
# Resposta Questão 2:
def produto_escalar(a, b):
    resultado = 0.0 if any(isinstance(x, float) for x in a + b) else 0
    i = 0
    while i < len(a):
        resultado += a[i] * b[i]
        i += 1
    return resultado

# Testes (Item b):
print(produto_escalar([1, 2, 3], [4, 5, 6]))        # Esperado: 32
print(produto_escalar([2.5, 0.0], [4.0, 10.0]))    # Esperado: 10.0
print(produto_escalar([], []))                    # Esperado: 0
\end{minted}
}

    % Questão 3
    \item Dada uma lista de elementos que representam os pontos de uma função $f(t)$ em um intervalo de tempo de interesse, separados por um passo $\Delta t$ fixo, uma forma de calcular a área sob a curva $f(t)$ neste intervalo é o método dos trapézios, dado por:
    \[
    \text{área} = \sum_{x=1}^{n} \frac{[f(x) + f(x-1)] \Delta t}{2}
    \]
    onde $x$ representa a posição na lista e $n$ é igual ao tamanho da lista menos 1.
    \begin{enumerate}[label=\alph*)]
        \item Escreva uma função em Python que recebe dois argumentos: uma lista de números que representam uma função $f(t)$ e o valor do passo $\Delta t$, e retorna a área calculada de acordo com a expressão acima.
        \item Teste a sua função de acordo com as entradas e saídas esperadas a seguir:
        \begin{itemize}
            \item \textbf{Entrada:} \code{area\_trapezio([2, 4, 6, 8], 1.0)} $\rightarrow$ \textbf{Saída Esperada:} \code{15.0}
            \item \textbf{Entrada:} \code{area\_trapezio([1, 1, 1, 1], 0.5)} $\rightarrow$ \textbf{Saída Esperada:} \code{1.5}
        \end{itemize}
    \end{enumerate}

    \vspace{0.5cm}
    \ifmostrarGabarito
\begin{minted}{python}
# Resposta Questão 3:
def area_trapezio(f, delta_t):
    if len(f) <= 1:
        return 0.0
    
    area = 0.0
    x = 1
    while x < len(f):
        area += ((f[x] + f[x - 1]) * delta_t) / 2.0
        x += 1
    return area

# Testes (Item b):
print(area_trapezio([2, 4, 6, 8], 1.0))  # Esperado: 15.0
print(area_trapezio([1, 1, 1, 1], 0.5))  # Esperado: 1.5
\end{minted}
\fi


    % Questão 4
    \item Escreva uma função em Python que calcula a diferença entre duas listas.
    \begin{enumerate}[label=\alph*)]
        \item A função deve receber duas listas por argumento e retornar uma nova lista que contém somente os números da primeira lista que não estão contidos na segunda lista. \textbf{Atenção:} Não utilize o operador \code{in}!
        \item Teste a sua função de acordo com as entradas e saídas esperadas a seguir:
        \begin{itemize}
            \item \textbf{Entrada:} \code{diferenca\_listas([1, 2, 3, 4, 5], [2, 4, 6])} $\rightarrow$ \textbf{Saída Esperada:} \code{[1, 3, 5]}
            \item \textbf{Entrada:} \code{diferenca\_listas([10, 20], [10, 20, 30])} $\rightarrow$ \textbf{Saída Esperada:} \code{[]}
            \item \textbf{Entrada:} \code{diferenca\_listas([1, 2, 3], [])} $\rightarrow$ \textbf{Saída Esperada:} \code{[1, 2, 3]}
        \end{itemize}
    \end{enumerate}
    \ifmostrarGabarito
\begin{minted}{python}
# Resposta Questão 4:
def diferenca_listas(lista1, lista2):
    resultado = []
    i = 0
    while i < len(lista1):
        elemento = lista1[i]
        # Verifica se 'elemento' está na lista2 sem usar o operador 'in'
        encontrado = False
        j = 0
        while j < len(lista2):
            if lista2[j] == elemento:
                encontrado = True
            j += 1
        
        if not encontrado:
            resultado.append(elemento)
        i += 1
    return resultado

# Testes (Item b):
print(diferenca_listas([1, 2, 3, 4, 5], [2, 4, 6])) # Esperado: [1, 3, 5]
print(diferenca_listas([10, 20], [10, 20, 30]))     # Esperado: []
print(diferenca_listas([1, 2, 3], []))               # Esperado: [1, 2, 3]
\end{minted}
\fi





    % Questão 5
    \item Escreva uma função em Python para contar a frequência de um caractere específico em uma string.
    \begin{enumerate}[label=\alph*)]
        \item A função deve receber por argumento uma string e um caractere (uma string de tamanho 1), e retornar o número de vezes que esse caractere aparece na string dada.
        \item Teste a sua função de acordo com as entradas e saídas esperadas a seguir:
        \begin{itemize}
            \item \textbf{Entrada:} \code{conta\_caractere("programacao", "a")} $\rightarrow$ \textbf{Saída Esperada:} \code{3}
            \item \textbf{Entrada:} \code{conta\_caractere("Computacao", "z")} $\rightarrow$ \textbf{Saída Esperada:} \code{0}
            \item \textbf{Entrada:} \code{conta\_caractere("", "a")} $\rightarrow$ \textbf{Saída Esperada:} \code{0}
        \end{itemize}
    \end{enumerate}

    \vspace{0.5cm}
    \ifmostrarGabarito
\begin{minted}{python}
# Resposta Questão 5:
def conta_caractere(texto, caractere):
    contador = 0
    i = 0
    while i < len(texto):
        if texto[i] == caractere:
            contador += 1
        i += 1
    return contador

# Testes (Item b):
print(conta_caractere("programacao", "a")) # Esperado: 3
print(conta_caractere("Computacao", "z"))  # Esperado: 0
print(conta_caractere("", "a"))            # Esperado: 0
\end{minted}
\fi


    % Questão 6
    \item Escreva uma função em Python que verifica se uma palavra ou frase é um palíndromo (lê-se da mesma forma de trás para frente e de frente para trás, ignorando maiúsculas, minúsculas e espaços).
    \begin{enumerate}[label=\alph*)]
        \item A função deve receber por argumento uma string e retornar \code{True} se a string for um palíndromo e \code{False} caso contrário. \textit{Dica:} você pode utilizar métodos de string para padronizar as letras e remover espaços.
        \item Teste a sua função de acordo com as entradas e saídas esperadas a seguir:
        \begin{itemize}
            \item \textbf{Entrada:} \code{eh\_palindromo("Arara")} $\rightarrow$ \textbf{Saída Esperada:} \code{True}
            \item \textbf{Entrada:} \code{eh\_palindromo("socorram me subi no onibus em marrocos")} $\rightarrow$ \textbf{Saída Esperada:} \code{True}
            \item \textbf{Entrada:} \code{eh\_palindromo("Python")} $\rightarrow$ \textbf{Saída Esperada:} \code{False}
        \end{itemize}
    \end{enumerate}
    

    \vspace{0.5cm}
    \ifmostrarGabarito{
\begin{minted}{python}
# Resposta Questão 6:
def eh_palindromo(texto):
    # Padroniza letras para minúsculas e remove espaços
    texto_limpo = texto.lower().replace(" ", "")
    
    # Verifica se a string é igual ao seu inverso
    return texto_limpo == texto_limpo[::-1]

# Testes (Item b):
print(eh_palindromo("Arara"))                                    # Esperado: True
print(eh_palindromo("socorram me subi no onibus em marrocos")) # Esperado: True
print(eh_palindromo("Python"))                                   # Esperado: False
\end{minted}
}


    % Questão 7
    \item Escreva uma função em Python para formatar e ocultar parte de um endereço de e-mail por motivos de privacidade.
    \begin{enumerate}[label=\alph*)]
        \item A função deve receber por argumento uma string contendo um e-mail válido (no formato \code{usuario@dominio.com}). Ela deve retornar o e-mail mascarado, onde apenas os 2 primeiros caracteres do nome do usuário ficam visíveis, substituindo o restante do nome do usuário por 4 asteriscos (\code{****}) mantendo o domínio intacto.
        \item Teste a sua função de acordo com as entradas e saídas esperadas a seguir:
        \begin{itemize}
            \item \textbf{Entrada:} \code{mascara\_email("joao@gmail.com")} $\rightarrow$ \textbf{Saída Esperada:} \\\code{"jo****@gmail.com"}
            \item \textbf{Entrada:} \code{mascara\_email("aluno@ufrj.br")} $\rightarrow$ \textbf{Saída Esperada:} \code{"al****@ufrj.br"}
        \end{itemize}
    \end{enumerate}

    \vspace{0.5cm}
    \ifmostrarGabarito
\begin{minted}{python}
# Resposta Questão 7:
def mascara_email(email):
    partes = email.split("@")
    usuario = partes[0]
    dominio = partes[1]
    
    usuario_mascarado = usuario[:2] + "****"
    return usuario_mascarado + "@" + dominio

# Testes (Item b):
print(mascara_email("joao.silva@gmail.com")) # Esperado: "jo****@gmail.com"
print(mascara_email("aluno@ufrj.br"))        # Esperado: "al****@ufrj.br"
\end{minted}
\fi


    % Questão 8
    \item Escreva uma função em Python para extrair apenas as palavras que possuem um comprimento mínimo a partir de uma frase.
    \begin{enumerate}[label=\alph*)]
        \item A função deve receber dois argumentos: uma string contendo uma frase e um número inteiro $k$. A função deve retornar uma lista de strings contendo apenas as palavras da frase cujo tamanho seja maior ou igual a $k$.
        \item Teste a sua função de acordo com as entradas e saídas esperadas a seguir:
        \begin{itemize}
            \item \textbf{Entrada:} \code{palavras\_longas("A estrutura de repeticao while e o tipo lista", 5)} $\rightarrow$ \textbf{Saída Esperada:} \code{['estrutura', 'repeticao', 'while', 'lista']}
            \item \textbf{Entrada:} \code{palavras\_longas("Python e uma linguagem incrivel", 7)} $\rightarrow$ \textbf{Saída Esperada:} \code{['linguagem', 'incrivel']}
            \item \textbf{Entrada:} \code{palavras\_longas("a b c", 2)} $\rightarrow$ \textbf{Saída Esperada:} \code{[]}
        \end{itemize}
    \end{enumerate}

    \ifmostrarGabarito{
\begin{minted}{python}
# Resposta Questão 8:
def palavras_longas(frase, k):
    palavras = frase.split()
    resultado = []
    i = 0
    while i < len(palavras):
        if len(palavras[i]) >= k:
            resultado.append(palavras[i])
        i += 1
    return resultado

# Testes (Item b):
print(palavras_longas("A estrutura de repeticao while e o tipo lista", 5))
# Esperado: ['estrutura', 'repeticao', 'while', 'lista']

print(palavras_longas("Python e uma linguagem incrivel", 7))
# Esperado: ['linguagem', 'incrivel']

print(palavras_longas("a b c", 2))
# Esperado: []
\end{minted}
}


    \item Escreva uma função em Python para sanitizar e extrair as palavras-chave (tags) de um texto formatado contendo dados brutos.
    \begin{enumerate}[label=\alph*)]
        \item A função deve receber por argumento uma string representando uma lista de tags brutas enviadas por um formulário (com espaçamentos extras no início/fim, letras maiúsculas e palavras separadas por hífen). A função deve realizar o tratamento do texto removendo os espaços das extremidades, convertendo todo o texto para minúsculas, substituindo os hífens (\code{-}) por espaços (\code{" "}) e, por fim, separando as tags por vírgula para retornar uma lista com as palavras-chave tratadas.
        
        \textit{Dica:} Para resolver esta questão, utilize os métodos de string \code{.strip()}, \code{.lower()}, \code{.replace()} e \code{.split()}.
        
        \item Teste a sua função de acordo com as entradas e saídas esperadas a seguir:
        \begin{itemize}
            \item \textbf{Entrada:} \code{processa\_tags("  PYTHON, MÁQUINA-DE-APRENDER, ESTRUTURA-DE-DADOS  ")} $\rightarrow$ \textbf{Saída Esperada:} \code{['python', 'máquina de aprender', 'estrutura de dados']}
            \item \textbf{Entrada:} \code{processa\_tags("  UFRJ, COMPUTAÇÃO-I, PROVA-FINAL  ")} $\rightarrow$ \textbf{Saída Esperada:} \code{['ufrj', 'computação i', 'prova final']}
            \item \textbf{Entrada:} \code{processa\_tags("  LÓGICA-DE-PROGRAMAÇÃO  ")} $\rightarrow$ \textbf{Saída Esperada:} 
            \\\code{['lógica de programação']}
        \end{itemize}
    \end{enumerate}
    \ifmostrarGabarito{
\begin{minted}{python}
# Resposta Questão 9:
def processa_tags(texto_bruto):
    # Utiliza .strip(), .lower(), .replace() e .split()
    texto_sanitizado = texto_bruto.strip().lower().replace("-", " ")
    tags = texto_sanitizado.split(",")
    return [tag.strip() for tag in tags]

# Testes (Item b):
print(processa_tags("  PYTHON, MÁQUINA-DE-APRENDER, ESTRUTURA-DE-DADOS  "))
# Esperado: ['python', 'máquina de aprender', 'estrutura de dados']

print(processa_tags("  UFRJ, COMPUTAÇÃO-I, PROVA-FINAL  "))
# Esperado: ['ufrj', 'computação i', 'prova final']

print(processa_tags("  LÓGICA-DE-PROGRAMAÇÃO  "))
# Esperado: ['lógica de programação']
\end{minted}
}


\end{enumerate}